\documentclass[11pt]{article}

\usepackage[margin=1in]{geometry}
\usepackage{amsmath,amssymb,amsthm,mathtools}
\usepackage{booktabs}
\usepackage{graphicx}
\usepackage{enumitem}
\usepackage{algorithm}
\usepackage{algpseudocode}
\usepackage{microtype}
\usepackage{xcolor}
\usepackage[hidelinks]{hyperref}

\newtheorem{definition}{Definition}[section]
\newtheorem{proposition}[definition]{Proposition}
\newtheorem{theorem}[definition]{Theorem}
\newtheorem{lemma}[definition]{Lemma}
\newtheorem{remark}[definition]{Remark}
\newtheorem{assumption}[definition]{Assumption}

\newcommand{\R}{\mathbb{R}}
\newcommand{\E}{\mathbb{E}}
\newcommand{\Pp}{\mathbb{P}}
\newcommand{\F}{\mathcal{F}}
\newcommand{\D}{\mathcal{D}}
\newcommand{\N}{\mathcal{N}}
\newcommand{\KL}{\operatorname{KL}}
\newcommand{\tr}{\operatorname{tr}}
\newcommand{\diag}{\operatorname{diag}}
\newcommand{\CA}{\mathrm{CA}}
\newcommand{\ind}{\mathbf{1}}

\title{Calibration-Aware Recurrent Neural Networks\\
for Multivariate Probabilistic Forecasting\\
under Distribution Shift}
\author{Jonathan Ma}
\date{Working paper\\\today}

\begin{document}

\hypersetup{pageanchor=false}
\maketitle

\begin{abstract}
Probabilistic forecasts are useful only when their stated uncertainty is credible.
Likelihood-trained recurrent neural networks may be miscalibrated, while Bayesian
model averaging is not guaranteed to remain calibrated when the likelihood, prior,
or conditional dynamics are misspecified. This paper develops the
Calibration-Aware Recurrent Neural Network (CA-RNN), adapting the
calibration-aware frequentist and Bayesian learning framework of Huang, Park, and
Simeone to continuous multivariate forecasting. The proposed training criterion
augments joint negative log likelihood with a differentiable discrepancy between
projected probability integral transforms (PITs) and the uniform distribution. A
sequential extension additionally penalizes lagged dependence in the projected
PIT sequence. The empirical design compares RNN, Bayesian RNN, CA-RNN, and
calibration-aware Bayesian RNN, and then extends the comparison to model
misspecification and distribution shift. In the four-asset application,
sequential CA-RNN reduces projected-PIT calibration error consistently across
optimization seeds and chronological market regimes, but does not uniformly
improve point or proper-score accuracy. The results identify a regime-dependent
calibration--accuracy tradeoff rather than universal dominance.
\end{abstract}

\tableofcontents
\clearpage
\hypersetup{pageanchor=true}

\section{Introduction}

\subsection{Motivation}

A point forecast reports one possible future; a probabilistic forecast reports a
distribution of possible futures. The latter is operationally meaningful only if
events assigned probability $\alpha$ occur at approximately rate $\alpha$, while
the distribution remains sufficiently sharp to discriminate among outcomes.
Modern neural forecasters can optimize likelihood successfully without achieving
reliable finite-sample calibration. Bayesian learning represents parameter
uncertainty, but its calibration guarantees depend on idealized assumptions about
the likelihood, prior, posterior computation, and data-generating mechanism.

Huang, Park, and Simeone propose calibration-aware Bayesian neural networks
(CA-BNNs), combining a data-dependent calibration regularizer with the
data-independent KL regularizer of variational Bayesian learning
\cite{huang2024cabl}. Their experiments compare frequentist, Bayesian,
calibration-aware frequentist, and calibration-aware Bayesian classifiers; vary
the calibration weight; and report expected calibration error and reliability
diagrams. Their conclusion identifies evaluation under distribution shift as a
direction for future work. This paper transfers that methodological template to
continuous multivariate sequential prediction and makes misspecification and
distribution shift a central part of the evaluation.

\subsection{Research questions}

The primary research question is:

\begin{quote}
Can calibration-aware training improve multivariate probabilistic RNN forecasts,
and does that improvement persist under model misspecification and distribution
shift without an unacceptable loss of sharpness or predictive accuracy?
\end{quote}

The supporting questions are:
\begin{enumerate}[label=RQ\arabic*.]
    \item How does the calibration weight affect projected calibration, proper
    predictive scores, sharpness, and point accuracy?
    \item Does calibration-aware Bayesian learning behave differently from its
    frequentist counterpart for dependent sequential data?
    \item Does penalizing serial dependence in projected PIT values improve
    sequential forecast adequacy beyond marginal PIT uniformity?
    \item How do the four model classes deteriorate under changes in volatility,
    dependence, tail behavior, conditional dynamics, and structural regime?
\end{enumerate}

\subsection{Contributions}

The contributions are methodological and empirical:
\begin{enumerate}
    \item a continuous-outcome analogue of calibration-aware neural training based
    on differentiable projected PIT discrepancies;
    \item an extension targeting temporal dependence in forecast PIT sequences;
    \item a variational Bayesian recurrent formulation combining calibration and
    KL regularization;
    \item a matched four-model experiment that parallels
    \cite{huang2024cabl}; and
    \item a controlled study of calibration degradation under misspecification and
    distribution shift, including a multivariate financial application.
\end{enumerate}

\subsection{Scope and non-claims}

Finite projected calibration is weaker than complete multivariate distributional
calibration. Calibration of a finite test set is also not a guarantee of
conditional calibration at every history. We therefore avoid claims
that a small collection of projections proves joint correctness. The proposed
regularizer is an explicit, diagnostically aligned training objective whose value
must be established empirically with proper scores and out-of-sample tests.

\section{Related Work}

\subsection{Calibration-aware Bayesian learning}

The organizing reference is \cite{huang2024cabl}. It distinguishes
data-dependent calibration regularization from the data-independent prior
regularization used by variational Bayesian inference. Its CA-FNN objective adds a
differentiable approximation of expected calibration error to cross-entropy, and
its CA-BNN objective places the same calibration penalty inside the variational
expectation while retaining a KL penalty. The paper uses Gaussian variational
parameters and the reparameterization trick, and its main comparison comprises
FNN, BNN, CA-FNN, and CA-BNN. The reported experiments emphasize calibration-weight
sweeps and reliability diagrams.

\subsection{Broader literature review to be completed}

This section is intentionally a structured stub. It will cover:
\begin{enumerate}
    \item probabilistic calibration, auto-calibration, and PIT diagnostics;
    \item proper scoring rules and calibration--sharpness decompositions;
    \item multivariate calibration, copula diagnostics, and projection methods;
    \item recurrent probabilistic forecasting and Bayesian recurrent networks;
    \item variational inference under misspecification;
    \item calibration under temporal dependence and distribution shift; and
    \item probabilistic forecasting and risk evaluation in financial time series.
\end{enumerate}

Only sources actually reviewed will be added to the bibliography. The local
literature archive is maintained in \texttt{docs/refs/}.

\section{Forecasting Setup and Calibration}

\subsection{Sequential prediction problem}

Let $(Y_t)_{t=1}^T$ be an adapted stochastic process with
$Y_t\in\R^d$ and filtration
$\F_t=\sigma(Y_1,\ldots,Y_t)$. Given a history length $H$, define
$X_t=(Y_{t-H+1},\ldots,Y_t)$. The task is to estimate a conditional predictive
distribution
\[
    P_t(\cdot)=\Pp(Y_{t+1}\in\cdot\mid\F_t)
\]
with a model $P_{\theta,t}$. Training, validation, and test blocks are ordered
chronologically. Standardization parameters are estimated on the training block
only.

\subsection{Joint Gaussian predictive head}

For the initial specification, the predictive law is
\begin{equation}
Y_{t+1}\mid\F_t,\theta
\sim \N_d\!\left(\mu_{\theta,t},\Sigma_{\theta,t}\right),
\qquad
\Sigma_{\theta,t}=L_{\theta,t}L_{\theta,t}^{\mathsf T}.
\label{eq:predictive}
\end{equation}
The network outputs the $d$ components of $\mu_{\theta,t}$ and the
$d(d+1)/2$ lower-triangular entries of $L_{\theta,t}$. An unconstrained diagonal
output $s_{t,j}$ is transformed as
\[
    L_{t,jj}=\operatorname{softplus}(s_{t,j})+\varepsilon,
    \qquad \varepsilon>0.
\]

\begin{proposition}[Validity of the predictive covariance]
For every network output, $\Sigma_{\theta,t}=L_{\theta,t}L_{\theta,t}^{\mathsf T}$
is symmetric positive definite.
\end{proposition}

\begin{proof}
Symmetry follows from transposition. Since every diagonal entry of the triangular
matrix $L_{\theta,t}$ is strictly positive, $L_{\theta,t}$ is nonsingular. Hence,
for any nonzero $x\in\R^d$,
\[
x^{\mathsf T}\Sigma_{\theta,t}x
=x^{\mathsf T}L_{\theta,t}L_{\theta,t}^{\mathsf T}x
=\|L_{\theta,t}^{\mathsf T}x\|_2^2>0.
\]
Thus $\Sigma_{\theta,t}$ is positive definite.
\end{proof}

\subsection{Negative log likelihood}

For an index set $I$, the mean joint Gaussian negative log likelihood is
\begin{align}
\mathcal L_{\mathrm{NLL}}(\theta;I)
&=-\frac{1}{|I|}\sum_{t\in I}
\log p_\theta(Y_{t+1}\mid X_t) \\
&=\frac{1}{2|I|}\sum_{t\in I}
\left[
d\log(2\pi)+2\sum_{j=1}^d\log L_{t,jj}
+e_t^{\mathsf T}\Sigma_{\theta,t}^{-1}e_t
\right],
\label{eq:nll}
\end{align}
where $e_t=Y_{t+1}-\mu_{\theta,t}$. Computation uses triangular solves rather
than an explicit matrix inverse.

\begin{proposition}[Strict propriety of the population log score]
Let $P_0(\cdot\mid\F_t)$ have density $p_0$ and let $p$ be any candidate
conditional density with common support. Then
\[
\E_{P_0}[-\log p(Y_{t+1}\mid\F_t)]
-\E_{P_0}[-\log p_0(Y_{t+1}\mid\F_t)]
=\E\!\left[\KL(P_0(\cdot\mid\F_t)\|P(\cdot\mid\F_t))\right]\geq0.
\]
Equality holds only when the conditional densities agree almost surely.
\end{proposition}

\begin{proof}
Condition on $\F_t$, subtract the two expected log scores, and recognize the
conditional KL divergence. Nonnegativity follows from Gibbs' inequality. Taking
expectations over $\F_t$ proves the result.
\end{proof}

This proposition explains why the calibration penalty must supplement rather than
replace a proper score: calibration alone can admit uninformative forecasts.

\subsection{Projected probability integral transforms}

Fix unit vectors $a_1,\ldots,a_R\in\R^d$. In the experiments they comprise the
$d$ coordinate directions, one equal-weight direction, and four random unit
directions fixed by a documented seed. Define
\begin{equation}
U_{t,r}=F_{\theta,t,r}(a_r^{\mathsf T}Y_{t+1}),
\label{eq:pit}
\end{equation}
where $F_{\theta,t,r}$ is the predictive CDF of $a_r^{\mathsf T}Y_{t+1}$.

\begin{proposition}[Projected Gaussian law]
Under \eqref{eq:predictive},
\[
a_r^{\mathsf T}Y_{t+1}\mid\F_t,\theta
\sim \N\!\left(a_r^{\mathsf T}\mu_{\theta,t},
a_r^{\mathsf T}\Sigma_{\theta,t}a_r\right),
\]
and therefore
\[
U_{t,r}=\Phi\!\left(
\frac{a_r^{\mathsf T}(Y_{t+1}-\mu_{\theta,t})}
{\sqrt{a_r^{\mathsf T}\Sigma_{\theta,t}a_r}}
\right).
\]
\end{proposition}

\begin{proof}
A linear transformation of a multivariate Gaussian is Gaussian. Its mean and
variance follow from linearity of expectation and the covariance transformation
rule. Applying its CDF gives the displayed PIT.
\end{proof}

\begin{theorem}[Sequential PIT property]
Suppose $F_{0,t,r}$ is the true, continuous conditional CDF of
$a_r^{\mathsf T}Y_{t+1}$ given $\F_t$. Then
$U_{t,r}=F_{0,t,r}(a_r^{\mathsf T}Y_{t+1})$ is uniformly distributed conditional
on $\F_t$. For a fixed $r$, the sequence $(U_{t,r})_t$ is independent and
$\operatorname{Uniform}(0,1)$.
\end{theorem}

\begin{proof}
For $u\in[0,1]$, continuity and the conditional probability integral transform
give
\[
\Pp(U_{t,r}\leq u\mid\F_t)=u.
\]
Thus $U_{t,r}$ is conditionally uniform and its conditional law does not depend on
$\F_t$. For times $t_1<\cdots<t_m$, iterated conditioning yields
\begin{align*}
\Pp(U_{t_1,r}\leq u_1,\ldots,U_{t_m,r}\leq u_m)
&=\E\left[
\ind_{\{U_{t_1,r}\leq u_1,\ldots,U_{t_{m-1},r}\leq u_{m-1}\}}
\Pp(U_{t_m,r}\leq u_m\mid\F_{t_m})
\right]\\
&=u_m\Pp(U_{t_1,r}\leq u_1,\ldots,U_{t_{m-1},r}\leq u_{m-1}).
\end{align*}
Induction gives $\prod_{j=1}^m u_j$, proving independence and uniformity.
\end{proof}

\begin{remark}
Uniformity for finitely many $a_r$ is necessary but not sufficient for equality of
multivariate distributions. If all one-dimensional projections agree, the
Cram\'er--Wold device identifies the joint distribution; a finite projection set
only provides a tractable approximation and diagnostic.
\end{remark}

Training uses the analytic Gaussian PIT above. Out-of-sample evaluation instead
uses the common sample-rank approximation
\[
 \widehat U_{t,r}=\frac{\sum_{s=1}^{S}
 \ind\{a_r^{\mathsf T}Y_{t+1}^{(s)}<a_r^{\mathsf T}Y_{t+1}\}+1/2}{S+1},
\]
with $S=500$ predictive draws. This supports the same diagnostic for deterministic
Gaussian forecasts and Bayesian posterior-predictive mixtures, at the cost of
finite-simulation error.

\section{Calibration-Aware RNN Methodology}

\subsection{Recurrent architecture}

For a single-layer GRU, with elementwise product $\odot$,
\begin{align*}
r_t&=\sigma(W_{ir}Y_t+b_{ir}+W_{hr}h_{t-1}+b_{hr}),\\
z_t&=\sigma(W_{iz}Y_t+b_{iz}+W_{hz}h_{t-1}+b_{hz}),\\
n_t&=\tanh(W_{in}Y_t+b_{in}+r_t\odot(W_{hn}h_{t-1}+b_{hn})),\\
h_t&=n_t+z_t\odot(h_{t-1}-n_t).
\end{align*}
Linear heads map the final hidden state to $\mu_{\theta,t}$ and the unconstrained
entries of $L_{\theta,t}$.

\subsection{Differentiable projected calibration loss}

Let $q_g=g/(G+1)$ for $g=1,\ldots,G$. The empirical CDF of PIT values is
nondifferentiable in model parameters because it uses indicators. Replace the
indicator $\ind\{U_{t,r}\leq q_g\}$ by
\[
s_\tau(q_g-U_{t,r})
=\sigma\!\left(\frac{q_g-U_{t,r}}{\tau}\right),\qquad \tau>0.
\]
For a consecutive training batch $B$, define
\begin{equation}
\mathcal L_{\mathrm{cal}}(\theta;B)
=\frac{1}{GR}\sum_{g=1}^G\sum_{r=1}^R
\left[
\frac{1}{|B|}\sum_{t\in B}s_\tau(q_g-U_{t,r})-q_g
\right]^2.
\label{eq:cal-loss}
\end{equation}
This is a grid approximation to an integrated squared CDF discrepancy, analogous
in purpose---but not in definition---to the differentiable AECE used for
classification in \cite{huang2024cabl}.

\begin{proposition}[Differentiability]
If the network outputs are differentiable in $\theta$ and
$a_r^{\mathsf T}\Sigma_{\theta,t}a_r>0$, then
$\mathcal L_{\mathrm{cal}}(\theta;B)$ is differentiable in $\theta$.
\end{proposition}

\begin{proof}
The Cholesky parameterization, affine projection, positive square root, Gaussian
CDF, logistic smoother, finite sums, and squaring operations are differentiable on
the stated domain. Their composition is therefore differentiable.
\end{proof}

\begin{proposition}[Zero-temperature limit]
Fix PIT values such that $U_{t,r}\neq q_g$ for every $t,r,g$. As
$\tau\downarrow0$, \eqref{eq:cal-loss} converges to the squared grid discrepancy
formed with the empirical PIT CDF.
\end{proposition}

\begin{proof}
For nonzero $x$, $\sigma(x/\tau)\to\ind\{x>0\}$ as $\tau\downarrow0$. Every sum
is finite, so the limit passes through the averages and squares. The exclusion of
ties removes the ambiguous boundary value.
\end{proof}

\subsection{Sequential calibration loss}

PIT uniformity does not rule out forecastable temporal structure. With
$\widetilde U_{t,r}=U_{t,r}-1/2$ and maximum lag $K$, define
\begin{equation}
\mathcal L_{\mathrm{seq}}(\theta;B)
=\frac{1}{RK}\sum_{r=1}^R\sum_{k=1}^K
\left[
\frac{1}{|B|-k}\sum_{t\in B:k\text{ available}}
\widetilde U_{t,r}\widetilde U_{t-k,r}
\right]^2.
\label{eq:seq-loss}
\end{equation}
Under the ideal conditional forecast, the population quantities targeted by
\eqref{eq:seq-loss} are zero by the sequential PIT theorem. The converse is not
true: finitely many zero autocovariances do not establish independence.

\subsection{Frequentist CA-RNN objective}

The proposed frequentist objective is
\begin{equation}
\widehat\theta_{\CA}
=\arg\min_\theta\left\{
\mathcal L_{\mathrm{NLL}}(\theta;\D)
+\lambda_{\mathrm{cal}}\mathcal L_{\mathrm{cal}}(\theta;\D)
+\lambda_{\mathrm{seq}}\mathcal L_{\mathrm{seq}}(\theta;\D)
\right\}.
\label{eq:ca-rnn}
\end{equation}
The replication experiment initially fixes $\lambda_{\mathrm{seq}}=0$, so that
\eqref{eq:ca-rnn} follows the source paper's task-loss-plus-calibration structure.
The sequential term is evaluated separately as a proposed extension.

\subsection{Bayesian CA-RNN objective}

Let every scalar neural parameter have variational posterior and prior
\[
q_\phi(\theta_j)=\N(\mu_j,\sigma_j^2),
\qquad
p(\theta_j)=\N(0,\sigma_0^2),
\qquad
\sigma_j=\operatorname{softplus}(\rho_j)+\epsilon.
\]
Sampling uses the reparameterization
\[
\theta_j=\mu_j+\sigma_j\varepsilon_j,
\qquad \varepsilon_j\sim\N(0,1).
\]
The independent-Gaussian KL term has closed form
\begin{equation}
\KL(q_\phi\|p)
=\sum_j\left[
\log\frac{\sigma_0}{\sigma_j}
+\frac{\sigma_j^2+\mu_j^2}{2\sigma_0^2}-\frac12
\right].
\label{eq:kl}
\end{equation}

The calibration-aware variational objective is
\begin{align}
\mathcal F^{\CA}(\phi)
={}&\E_{\theta\sim q_\phi}\left[
\mathcal L_{\mathrm{NLL}}(\theta;\D)
+\lambda_{\mathrm{cal}}\mathcal L_{\mathrm{cal}}(\theta;\D)
+\lambda_{\mathrm{seq}}\mathcal L_{\mathrm{seq}}(\theta;\D)
\right]\\
&+\frac{\beta}{N}\KL(q_\phi(\theta)\|p(\theta)).
\label{eq:ca-brnn}
\end{align}
This is the sequential continuous-outcome analogue of the CA-BNN free energy in
\cite{huang2024cabl}. The factor $1/N$ matches the mean per-observation likelihood
convention. The weight $\beta$ remains a validation-selected generalized-Bayes
hyperparameter.

\begin{proposition}[Unbiased reparameterized gradient for the expectation]
Assume differentiation and expectation may be interchanged. If
$\theta^{(m)}=\mu+\sigma(\rho)\odot\varepsilon^{(m)}$ with independent
$\varepsilon^{(m)}\sim\N(0,I)$, then the Monte Carlo gradient of the expected
data-dependent part of \eqref{eq:ca-brnn} is unbiased.
\end{proposition}

\begin{proof}
Write the expectation with respect to the parameter-free distribution of
$\varepsilon$. Under the interchange assumption,
\[
\nabla_\phi\E_\varepsilon[f(\mu+\sigma(\rho)\odot\varepsilon)]
=\E_\varepsilon\nabla_\phi f(\mu+\sigma(\rho)\odot\varepsilon).
\]
The sample average of independent terms has this expectation. The KL gradient is
computed analytically from \eqref{eq:kl}.
\end{proof}

\begin{remark}[Dependent mini-batches]
Unlike independent classification observations, overlapping time-series windows
are dependent. Consecutive batches are required to evaluate
$\mathcal L_{\mathrm{seq}}$, but their nonlinear calibration loss is not asserted
to be an unbiased estimate of the full-sample calibration loss. Batch length and
history overlap are therefore methodological hyperparameters and sensitivity
targets.
\end{remark}

In the implementation, both calibration penalties are batch-local: lag pairs
crossing mini-batch boundaries are omitted. Epoch summaries average batch means,
so the final shorter batch receives the same summary weight as a full batch. The
displayed dataset-level objectives are the target criteria; stochastic
optimization uses this documented batch-local surrogate.

\subsection{Training algorithms}

\begin{algorithm}[ht]
\caption{Frequentist CA-RNN training}
\label{alg:ca-rnn}
\begin{algorithmic}[1]
\Require Ordered observations $\D$; history $H$; projections $A$; weights
$\lambda_{\mathrm{cal}},\lambda_{\mathrm{seq}}$; epoch limit $E$;
gradient threshold $c$
\Ensure Fitted parameters $\widehat\theta$ selected without test information
\State Split $\D$ chronologically into $\D_{\mathrm{tr}},\D_{\mathrm{val}},
\D_{\mathrm{te}}$
\State Fit standardization on $\D_{\mathrm{tr}}$ and construct length-$H$ windows
\State Initialize $\theta$ and fix $A=(a_1,\ldots,a_R)^{\mathsf T}$
\For{$e=1,\ldots,E$}
    \For{consecutive mini-batch $B\subset\D_{\mathrm{tr}}$}
        \State $(\mu_{\theta,t},L_{\theta,t})_{t\in B}
        \gets \Call{CARNNForward}{B,\theta}$
        \State Compute $\mathcal L_{\mathrm{NLL}}$ by \eqref{eq:nll} and
        projected PITs by \eqref{eq:pit}
        \State Compute $\mathcal L_{\mathrm{cal}}$ by \eqref{eq:cal-loss} and
        $\mathcal L_{\mathrm{seq}}$ by \eqref{eq:seq-loss}
        \State $\mathcal J\gets\mathcal L_{\mathrm{NLL}}
        +\lambda_{\mathrm{cal}}\mathcal L_{\mathrm{cal}}
        +\lambda_{\mathrm{seq}}\mathcal L_{\mathrm{seq}}$
        \State $g\gets\Call{ClipGradNorm}{\nabla_\theta\mathcal J,c}$
        \State Update $\theta$ with AdamW using $g$
    \EndFor
    \State Evaluate the prespecified selection score on $\D_{\mathrm{val}}$
    \State Save $\theta$ if the validation score improves
    \If{the patience criterion is met} \State \textbf{break} \EndIf
\EndFor
\State Restore the selected checkpoint $\widehat\theta$
\State Evaluate $\widehat\theta$ once on $\D_{\mathrm{te}}$
\end{algorithmic}
\end{algorithm}

For fixed-budget causal loss ablations, checkpoint selection and the patience test
in Algorithm~\ref{alg:ca-rnn} are disabled, and all specifications are compared
after the same number of epochs.

\begin{algorithm}[ht]
\caption{Variational calibration-aware Bayesian RNN training}
\label{alg:ca-brnn}
\begin{algorithmic}[1]
\Require Ordered observations $\D$; projections $A$; prior scale $\sigma_0$;
Monte Carlo count $M$; loss weights $\lambda_{\mathrm{cal}},
\lambda_{\mathrm{seq}},\beta$; epoch limit $E$
\Ensure Variational parameters $\widehat\phi=(\widehat\mu,\widehat\rho)$
\State Construct chronological splits and training-only standardization as in
Algorithm~\ref{alg:ca-rnn}
\State Initialize $\phi=(\mu,\rho)$ and set
$\sigma(\rho)=\operatorname{softplus}(\rho)+\epsilon$
\For{$e=1,\ldots,E$}
    \For{consecutive mini-batch $B\subset\D_{\mathrm{tr}}$}
        \For{$m=1,\ldots,M$}
            \State Draw $\varepsilon^{(m)}\sim\N(0,I)$ and set
            $\theta^{(m)}\gets\mu+\sigma(\rho)\odot\varepsilon^{(m)}$
            \State Use $\theta^{(m)}$ consistently at every recurrent time step
            \State Compute $\mathcal J^{(m)}_{\mathrm{data}}\gets
            \mathcal L_{\mathrm{NLL}}^{(m)}+\lambda_{\mathrm{cal}}
            \mathcal L_{\mathrm{cal}}^{(m)}+\lambda_{\mathrm{seq}}
            \mathcal L_{\mathrm{seq}}^{(m)}$
        \EndFor
        \State $\widehat{\mathcal F}^{\CA}\gets M^{-1}\sum_{m=1}^M
        \mathcal J^{(m)}_{\mathrm{data}}+(\beta/N)\KL(q_\phi\Vert p)$
        \State Differentiate $\widehat{\mathcal F}^{\CA}$, clip its gradient,
        and update $\phi$ with AdamW
    \EndFor
    \State Evaluate validation NLL at posterior-mean weights and update the
    selected checkpoint
    \If{the patience criterion is met} \State \textbf{break} \EndIf
\EndFor
\State Restore $\widehat\phi$
\For{each test origin $t$ and predictive draw $s$}
    \State Draw $\theta^{(s)}\sim q_{\widehat\phi}$ and then
    $Y_{t+1}^{(s)}\sim p_{\theta^{(s)}}(\cdot\mid\F_t)$
\EndFor
\State Evaluate the posterior predictive mixture on $\D_{\mathrm{te}}$
\end{algorithmic}
\end{algorithm}

Algorithm~\ref{alg:ca-brnn} is the recurrent continuous-outcome adaptation of the
CA-BNN procedure in \cite{huang2024cabl}.

\section{Model Misspecification and Distribution Shift}

\subsection{Distinguishing misspecification from shift}

Model misspecification occurs when the true conditional law $P_0(Y_{t+1}\mid
\F_t)$ is outside the fitted family, even if the data law is stationary.
Distribution shift occurs when the evaluation law differs from the development
law. The two can coexist: a Gaussian predictive head is misspecified under
heavy-tailed innovations, while a transition from Gaussian training innovations
to heavy-tailed test innovations is also a distribution shift.

Let $P_0$ denote the reference law and let $P_{\delta,s}$ denote stress family $s$
at severity $\delta\geq0$, with $P_{0,s}=P_0$. For metric $M$, define degradation
\begin{equation}
\Delta_M^{(m,s)}(\delta)
=M(P_{\delta,s},\widehat P_m)-M(P_{0,s},\widehat P_m),
\label{eq:degradation}
\end{equation}
where $m$ indexes the forecasting method. For metrics in which smaller is better,
positive $\Delta_M$ denotes deterioration. Coverage is reported as absolute error
from the nominal level before differencing.

\subsection{Stress families}

The following are prospective isolated stress families for a confirmatory study;
they were not all executed factorially in the present experiments:
\begin{enumerate}
    \item \textbf{Scale shift:} innovation covariance is multiplied by a severity
    function while dependence is held fixed.
    \item \textbf{Dependence shift:} correlations or a dynamic correlation state
    change while marginal variances are controlled.
    \item \textbf{Tail shift:} Gaussian innovations transition to Student-$t$,
    mixtures, or asymmetric downside contamination.
    \item \textbf{Volatility misspecification:} ARCH/GARCH or stochastic-volatility
    dynamics violate the Gaussian head's conditionally homoskedastic approximation.
    \item \textbf{Mean-dynamics shift:} nonlinear interactions or an abrupt change
    in recurrent coefficients alter the conditional location.
\end{enumerate}

The implemented developmental curve uses common random numbers and jointly changes
scale and common correlation. Models are selected without examining its stressed
test metrics and are frozen after fitting at $\delta=0$. Tail, volatility, and
mean-dynamics regimes appear in separate developmental simulations, not in a
completed isolated-factor robustness experiment.

\subsection{Robustness estimands}

For paired realization $b$, method $m$, comparator $c$, and metric $M$, define
\[
D_b^{m,c}(\delta)=M_b^m(\delta)-M_b^c(\delta).
\]
For a future experiment with $B$ independent DGP realizations, the intended
summary is
\[
\overline D^{m,c}(\delta)=\frac1B\sum_{b=1}^B D_b^{m,c}(\delta),
\qquad
\operatorname{MCSE}(\overline D)=
\frac{s_D(\delta)}{\sqrt B}.
\]
Neural-seed repetitions on a prespecified subset separate DGP variability from
optimization variability. Failure frequencies and runtimes are part of the
estimand rather than discarded implementation details.

\section{Experimental Design}

\subsection{Four-model replication matrix}

The core experiment follows the comparison logic of
\cite{huang2024cabl}:
\begin{table}[ht]
\centering
\caption{Mapping from calibration-aware classification to forecasting.}
\begin{tabular}{lll}
\toprule
Source-paper model & Forecasting model & Objective components \\
\midrule
FNN & RNN & NLL \\
BNN & BRNN & expected NLL $+$ KL \\
CA-FNN & CA-RNN & NLL $+$ projected-PIT calibration \\
CA-BNN & CA-BRNN & expected NLL/calibration $+$ KL \\
\bottomrule
\end{tabular}
\end{table}

Architectures, histories, splits, optimizers, stopping rules, and predictive heads
are matched wherever the Bayesian formulation permits. The same projection matrix
is used for all models in a comparison.

\subsection{Calibration-weight experiment}

Analogously to the source paper's plot of ECE and accuracy against $\lambda$, this
study plots projected-PIT calibration error and predictive performance against
$\lambda_{\mathrm{cal}}$. The paired panels will contain:
\begin{enumerate}
    \item projected calibration discrepancy and reliability curves;
    \item Energy Score, NLL where comparable, RMSE, interval score, and width.
\end{enumerate}
The result of interest is a Pareto relationship, not calibration alone. Separate
curves are required for CA-RNN and CA-BRNN.

\subsection{Sequential ablation}

The minimum ablation is:
\begin{table}[ht]
\centering
\caption{Loss-component ablation.}
\begin{tabular}{lccc}
\toprule
Specification & NLL & $\mathcal L_{\mathrm{cal}}$ & $\mathcal L_{\mathrm{seq}}$ \\
\midrule
RNN & \checkmark & & \\
CA-RNN & \checkmark & \checkmark & \\
CA-RNN (sequential) & \checkmark & \checkmark & \checkmark \\
\bottomrule
\end{tabular}
\end{table}
The Bayesian versions repeat this comparison with the KL term always present.
Additional ablations vary coordinate-only versus augmented projections, smoothing
temperature, batch length, lag order, prior scale, and $\beta$.

To prevent checkpoint selection from masquerading as a calibration effect, the
primary loss-component ablation uses identical initial parameters, chronological
mini-batches, optimization budgets, and predictive random numbers, with early
stopping disabled. A secondary operational comparison may use early stopping, but
all methods must then be selected by the same validation NLL rule.

\subsection{Controlled simulations}

Development simulations include Gaussian, heavy-tailed, heteroskedastic, and
nonlinear regimes. A future confirmatory experiment should use independently
generated data for every regime, sample size, shift severity, and DGP seed, with
the following fixed in advance:
\begin{itemize}
    \item sample sizes and number of independent realizations;
    \item neural-seed subset and number of optimization repetitions;
    \item calibration and KL grids;
    \item projection directions and smoothing grid;
    \item primary metrics and failure definitions; and
    \item rules for early stopping and invalid numerical runs.
\end{itemize}

\subsection{Financial application}

The application uses adjusted daily prices for AAPL, JPM, XOM, and WMT beginning
in 2007 and converts them to aligned log returns. The four assets provide a small,
computationally manageable multivariate system with heterogeneous sector exposure.
All splits are chronological. Development selected
$\lambda_{\mathrm{cal}}=100$ and $\lambda_{\mathrm{seq}}=2000$ before the
multi-seed and rolling-regime analyses. The focused frequentist comparison uses a
32-unit GRU, a 20-day history, a joint Gaussian head, 500 forecast draws, and
validation-NLL early stopping.

Standardization is fitted on each training partition. Validation contexts carry
the final 20 training observations, and test contexts carry the final 20
validation observations; every evaluation target therefore retains an available
historical context without using future information. RMSE, Energy Score, interval
width, and interval score are reported in training-standardized units. Coverage
and PIT diagnostics are invariant to the componentwise affine transformation.
The two penalty weights were chosen after inspecting developmental simulation
and ablation outcomes, rather than by a pristine validation-only or preregistered
selection procedure.

The rolling analysis uses expanding estimation samples and four non-overlapping
test periods: 2017--2019 (pre-pandemic), 2020--2021 (pandemic), 2022--2023
(inflation and monetary tightening), and 2024--2025 (recent). Each window has a
strictly subsequent validation period and five matched optimization seeds. These
labels describe calendar periods rather than uniquely identified causal regimes.
Because the full panel was observed during development and the final two windows
overlap the initial test period, this is temporal robustness evidence rather than
a pristine preregistered confirmatory experiment.

\subsection{Paired uncertainty design}

For model $m$, comparator $c$, seed $s$, and forecast origin $t$, let
$d_{s,t}^{m,c}$ be a paired additive-score difference. The estimand is
\[
 \widehat\Delta^{m,c}=\frac{1}{S}\sum_{s=1}^S\frac{1}{T}
 \sum_{t=1}^T d_{s,t}^{m,c}.
\]
The hierarchical bootstrap resamples optimization seeds and, within each selected
seed, matched forecast origins in moving blocks of 20 trading days. Projected-PIT
calibration, coverage error, and PIT dependence are recomputed in every bootstrap
draw. The initial financial analysis uses 2,000 draws and each rolling regime uses
1,000. Raw coverage and width remain descriptive; the inferential coverage loss is
\[
 E_{\mathrm{cov}}=|\widehat C_{0.90}-0.90|.
\]
Negative candidate-minus-comparator differences favor the candidate for loss
metrics. An interval containing zero is called unresolved, not evidence of
equality. Inference conditions on the architecture, tuning path, projection set,
assets, and observed calendar windows.

\subsection{Metrics and figures}

The tables report mean performance, paired differences, and uncertainty
intervals. The diagnostic figures include:
\begin{enumerate}
    \item calibration and predictive metrics versus $\lambda_{\mathrm{cal}}$;
    \item nominal-versus-empirical coverage curves;
    \item projected PIT histograms or CDF reliability curves;
    \item PIT autocorrelation diagnostics;
    \item degradation curves versus shift severity; and
    \item a calibration--sharpness or calibration--Energy Score frontier.
\end{enumerate}

\section{Results}

\subsection{Implementation validation}

The implementation tests cover tensor dimensions, covariance
positive-definiteness, loss evaluation, forecasting, chronological data handling,
and paired-bootstrap identities. Eighteen automated tests pass. The notebooks also
completed the simulation, shift, financial, optimization-seed, and rolling-regime
workflows. Successful execution establishes implementation consistency but is not
itself statistical validation.

\subsection{Four-model comparison}

The initial matched comparison showed that calibration-aware objectives alter the
forecast distribution but do not yield uniform dominance. CA-RNN made only a
small projected-calibration improvement over RNN, while CA-BRNN did not improve
on BRNN. Bayesian averaging did not remove the calibration--score tradeoff. The
frequentist RNN, CA-RNN, and sequential
CA-RNN therefore form the focused empirical comparison; BRNN and CA-BRNN remain
supporting replications of the source paper's model matrix.

\subsection{Calibration-weight sensitivity}

The exploratory sweep over
$\lambda_{\mathrm{cal}}\in\{0,1,10,50,100,250,500\}$ displayed a Pareto
tradeoff: projected calibration improved through the intermediate penalty range,
whereas likelihood and proper-score performance deteriorated. The value 100 was
retained as a developmental compromise rather than selected to maximize a final
test result. Calibration improvement alone is therefore not interpreted as a
better probabilistic forecast.

\subsection{Sequential calibration ablation}

Relative to marginal CA-RNN, the sequential specification consistently repaired
predictive performance. In the ten-seed financial experiment it improved MSE,
Energy Score, interval score, nominal-coverage error, and projected-PIT
calibration for every seed. It widened intervals, but simultaneous coverage and
interval-score improvements show that the result is not merely unpenalized
widening. Improvements in PIT temporal dependence were smaller and generally
unresolved. The sequential penalty is empirically important, but the evidence
does not establish PIT independence.

\subsection{Misspecification and distribution shift}

The frozen-model experiment jointly increased innovation scale and common
correlation from severity zero to one. All specifications deteriorated. At zero,
RNN had Energy Score 1.2094, interval score 3.8581, and projected calibration
error 0.001068, compared with 1.2124, 3.9251, and 0.001545 for sequential CA-RNN.
At maximum severity, the ordering reversed: RNN obtained 2.3671, 10.3070, and
0.009370, while sequential CA-RNN obtained 2.3416, 10.0120, and 0.007728. This
single-seed combined stress is developmental evidence rather than an isolated
causal test of scale or dependence. Absolute levels and paired degradation are
considered together because a worse nominal starting point can mechanically
reduce measured deterioration.

\subsection{Financial results}

Table~\ref{tab:seed-results} reports averages over ten matched optimization seeds
on the original financial test period. Sequential CA-RNN reduced projected-PIT
calibration error relative to RNN for all ten seeds, from $0.001382$ to
$0.001185$, or approximately 14\%. Its Energy Score and RMSE were worse for all
ten seeds. The interval-score difference was unresolved, with a 95\% hierarchical
interval of $[-0.0047,0.0235]$, as was the nominal-coverage error difference.

\begin{table}[ht]
\centering
\caption{Mean financial performance over ten matched optimization seeds. Lower
is better except for coverage, whose nominal target is 0.90.}
\label{tab:seed-results}
\begin{tabular}{lrrrrr}
\toprule
Model & RMSE & Energy & Interval & Coverage & PIT calibration \\
\midrule
RNN & 0.9368 & 1.1394 & 4.0509 & 0.8725 & 0.001382 \\
CA-RNN & 0.9474 & 1.1510 & 4.0979 & 0.8650 & 0.001265 \\
Sequential CA-RNN & 0.9448 & 1.1478 & 4.0612 & 0.8748 & 0.001185 \\
\bottomrule
\end{tabular}
\end{table}

Plain CA-RNN isolates why the extension matters. Against RNN it had worse MSE,
Energy Score, interval score, and coverage error despite lower projected-PIT
error. Sequential CA-RNN then improved MSE by $0.00483$, Energy Score by
$0.00329$, interval score by $0.03673$, and coverage error by $0.00979$ relative
to CA-RNN; every corresponding hierarchical interval excluded zero. It repairs a
substantial part, but not all, of the cost of marginal calibration.

\subsection{Rolling market-regime results}

Table~\ref{tab:regime-results} gives five-seed means. Sequential CA-RNN had lower
projected-PIT calibration error than RNN in all four regimes. Hierarchical
intervals favored it in the pre-pandemic, inflation/tightening, and recent
windows; the pandemic difference was unresolved. In contrast, its Energy Score
was lower only pre-pandemic, where the difference was unresolved.

\begin{table}[ht]
\centering
\caption{Five-seed regime means for RNN (R) and sequential CA-RNN (S).}
\label{tab:regime-results}
\begin{tabular}{lrrrrrr}
\toprule
& \multicolumn{2}{c}{Energy Score} & \multicolumn{2}{c}{Interval score}
& \multicolumn{2}{c}{PIT calibration} \\
Regime & R & S & R & S & R & S \\
\midrule
Pre-pandemic & 0.8629 & 0.8626 & 3.1123 & 3.1027 & 0.001914 & 0.001503 \\
Pandemic & 1.5016 & 1.5029 & 5.7371 & 5.6629 & 0.002995 & 0.002883 \\
Inflation/tightening & 1.2809 & 1.2884 & 4.5862 & 4.5396 & 0.001842 & 0.001465 \\
Recent & 1.0237 & 1.0289 & 3.6648 & 3.6935 & 0.001224 & 0.000904 \\
\bottomrule
\end{tabular}
\end{table}

The interval-score ordering is regime-dependent. Sequential CA-RNN significantly
outperformed RNN during the pandemic and inflation/tightening periods, was
unresolved pre-pandemic, and significantly underperformed recently. Coverage
error similarly favored sequential CA-RNN pre-pandemic and during
inflation/tightening, was unresolved during the pandemic, and favored RNN
recently. These reversals directly answer the distribution-shift question.

Across regimes, sequential CA-RNN improved Energy Score and interval score over
marginal CA-RNN directionally in four of four windows. These improvements were
resolved in the pandemic, inflation/tightening, and recent windows, while the
pre-pandemic differences were unresolved. Projected calibration improved over
marginal CA-RNN in three of four windows.

\begin{figure}[ht]
\centering
\IfFileExists{../results/market_regimes/regime_metric_profiles.png}{%
\includegraphics[width=\textwidth]{../results/market_regimes/regime_metric_profiles.png}%
}{\fbox{\parbox{0.9\textwidth}{Run notebook 09 to generate the rolling-regime
metric profile.}}}
\caption{Mean Energy Score, interval score, and projected-PIT calibration error
across five matched seeds and four chronological test regimes. Lower is better.}
\label{fig:regime-profiles}
\end{figure}

\subsection{Interval-score mechanism under temporal shift}

To explain the regime reversal, the retained forecasts were analyzed without
refitting. For a 90\% interval $[L_{t,j},U_{t,j}]$, the componentwise interval
score decomposes exactly as
\begin{equation}
 S_{t,j}=\underbrace{U_{t,j}-L_{t,j}}_{\text{width}}
 +\underbrace{20(L_{t,j}-Y_{t,j})\ind\{Y_{t,j}<L_{t,j}\}}_{\text{lower-tail penalty}}
 +\underbrace{20(Y_{t,j}-U_{t,j})\ind\{Y_{t,j}>U_{t,j}\}}_{\text{upper-tail penalty}}.
 \label{eq:interval-decomposition}
\end{equation}
Table~\ref{tab:interval-mechanism} reports the mean sequential-CA-RNN-minus-RNN
difference in each component. The arithmetic identity in
\eqref{eq:interval-decomposition} means the final column is exactly the sum of the
preceding three columns.

\begin{table}[ht]
\centering
\caption{Decomposition of the interval-score difference between sequential
CA-RNN and RNN. Negative total differences favor sequential CA-RNN.}
\label{tab:interval-mechanism}
\begin{tabular}{lrrrr}
\toprule
Regime & Width & Lower penalty & Upper penalty & Total \\
\midrule
Pre-pandemic & $-0.0439$ & $+0.0241$ & $+0.0103$ & $-0.0096$ \\
Pandemic & $+0.0420$ & $-0.0599$ & $-0.0563$ & $-0.0742$ \\
Inflation/tightening & $+0.2342$ & $-0.1709$ & $-0.1098$ & $-0.0466$ \\
Recent & $-0.0185$ & $+0.0268$ & $+0.0205$ & $+0.0287$ \\
\bottomrule
\end{tabular}
\end{table}

The pandemic and inflation/tightening improvements arise through beneficial
widening: the added width is more than offset by reductions in both tail-miss
penalties. In the recent period, sequential CA-RNN instead produces slightly
narrower intervals and larger penalties in both tails. Before the pandemic, the
mechanism is different again: reduced width outweighs the additional tail
penalties, yielding a small and statistically unresolved sharpness benefit.

The stress-period effect is concentrated in forecasts preceded by elevated
volatility. Sequential CA-RNN's interval-score differences in the pandemic are
$+0.0198$, $-0.0725$, and $-0.1705$ for low-, medium-, and high-volatility
origins, respectively. During inflation/tightening they are $+0.0209$, $-0.0562$,
and $-0.1045$. By contrast, the recent-period differences are positive in every
state: $+0.0318$, $+0.0238$, and $+0.0305$. Volatility states use trailing 20-day
realized volatility known before each forecast and within-regime terciles.

The improvements are also asset-concentrated. XOM accounts for most of the
inflation/tightening gain, with an interval-score difference of $-0.2364$;
AAPL, JPM, and WMT have differences of $-0.0051$, $+0.0221$, and $+0.0331$.
Pandemic improvements occur for AAPL, XOM, and WMT but not JPM. Recent losses are
concentrated in AAPL, XOM, and WMT, while JPM improves slightly. These results
preclude interpreting an aggregate gain as uniform across assets.

The train-to-test descriptors support, but do not prove, a scale-and-tail
interpretation. Pandemic volatility is 1.26 times its training level and the mean
standardized fourth moment is 29.87. The corresponding values are 1.00 and 12.08
during inflation/tightening, versus 0.83 and 6.37 recently. The recent period has
the largest standardized location shift and correlation shift of the four
windows. Thus sequential calibration appears most useful for persistent scale and
tail errors, while lower-volatility location or dependence changes remain harder.

Importantly, projected and coordinate PIT calibration still improve recently even
as 90\% interval performance deteriorates. The calibration loss averages 19 CDF
grid levels and multiple projections; it does not target the 5\% and 95\% quantiles
specifically. Better global PIT uniformity can therefore coexist with worse tail
coverage. The exact score decomposition is established by the retained forecasts,
whereas the proposed economic mechanism is post-hoc and requires new data for
confirmation.

\begin{figure}[ht]
\centering
\IfFileExists{../results/interval_mechanisms/interval_score_decomposition.png}{%
\includegraphics[width=0.82\textwidth]{../results/interval_mechanisms/interval_score_decomposition.png}%
}{\fbox{\parbox{0.8\textwidth}{Run notebook 10 to generate the interval-score
decomposition.}}}
\caption{Exact decomposition of sequential-CA-RNN-minus-RNN interval-score
differences. Tail reductions offset widening in the pandemic and
inflation/tightening periods, but not in the recent period.}
\label{fig:interval-mechanism}
\end{figure}

\section{Discussion}

\subsection{Interpretation}

Three conclusions follow. First, a marginal calibration penalty is insufficient
for dependent multivariate forecasting: it can improve its targeted diagnostic
while harming proper scores and nominal coverage. Second, the sequential loss is
not cosmetic; it consistently improves marginal CA-RNN and recovers much of its
predictive-performance loss. Third, sequential CA-RNN is a calibration-oriented
method rather than a uniformly more accurate forecaster. Its projected-calibration
advantage is stable across seeds and calendar periods, whereas its proper-score
advantage depends on regime and scoring rule.

The disagreement between Energy Score and interval score is substantively useful.
Energy Score assesses the multivariate predictive distribution, whereas the
marginal interval score rewards coverage and sharpness asset by asset. During the
pandemic and inflation/tightening periods, sequential CA-RNN improves the latter
without improving the former. Its practical value therefore depends on whether an
application prioritizes marginal risk intervals or overall multivariate accuracy.

\subsection{Limitations}

Limitations include finite projection sets, four assets, one-step horizons,
sensitivity to smoothing and batching, mean-field variational misspecification, a
Gaussian predictive head, and limited ability to establish conditional
calibration from finite dependent sequences. Regime labels are calendar-based and
partly specified after broader exposure to the data. The hierarchical bootstrap
represents seed and within-window sampling variation but conditions on the
architecture, tuning path, assets, projections, and chosen windows. Moving-block
resampling can distort PIT autocorrelation at artificial block boundaries, so
temporal-dependence intervals are interpreted cautiously. Finally, improved
finite-projection calibration does not imply full joint or conditional
calibration.

\subsection{Future work}

The most direct methodological extension is a tail-aware sequential calibration
objective. The current loss weights all 19 PIT grid levels and all fixed
projections equally. A weighted version could emphasize application-relevant tail
levels,
\begin{equation}
 \mathcal L_{\mathrm{tail}}
 =\frac{1}{R\sum_g w_g}\sum_{r=1}^R\sum_{g=1}^G w_g
 \left[\frac{1}{|B|}\sum_{t\in B}s_\tau(q_g-U_{t,r})-q_g\right]^2,
 \label{eq:tail-calibration}
\end{equation}
with larger prespecified $w_g$ near $q=0.05$ and $q=0.95$. A complementary
extension would penalize lagged dependence in tail exceedance indicators rather
than only covariance of centered PITs. Either extension must be tuned on
validation data and evaluated with global calibration and proper scores to avoid
improving one nominal interval at the expense of the rest of the distribution.

Further work should test whether the observed high-volatility benefit replicates
for new assets, frequencies, forecast horizons, and genuinely held-out calendar
periods. Conditional calibration diagnostics could use volatility, market return,
and liquidity states fixed before evaluation. Heavy-tailed or mixture predictive
heads may separate deficiencies in the Gaussian output law from deficiencies in
the calibration objective. Finally, uncertainty for PIT dependence should use a
procedure that avoids correlations introduced at moving-block boundaries.

\section{Conclusion}

This paper adapts calibration-aware Bayesian learning to continuous multivariate
recurrent forecasting and introduces a temporal calibration penalty based on
lagged projected-PIT covariance. Calibration-aware training reliably changes
forecast calibration, but marginal regularization alone incurs predictive costs.
The sequential extension consistently improves the marginal CA-RNN specification
and produces lower projected-PIT calibration error than a likelihood-trained RNN
across ten optimization seeds and four chronological market regimes.

The method is not uniformly superior. RNN retains better point and Energy Score
performance in much of the financial evidence, and interval-score advantages
reverse across regimes. The contribution is therefore not a dominance claim. It
is evidence that temporal calibration must be addressed explicitly, that doing so
repairs the direct calibration-aware adaptation, and that the remaining
calibration--accuracy tradeoff varies systematically under temporal distribution
shift.

\appendix

\section{Implementation Correspondence}

\begin{table}[ht]
\centering
\caption{Mathematical objects and current source modules.}
\small
\begin{tabular}{ll}
\toprule
Object & Source location \\
\midrule
GRU and joint Gaussian head & \texttt{src/innovcal/ca\_rnn/model.py} \\
Mean-field variational layers & \texttt{src/innovcal/ca\_rnn/layers.py} \\
Cholesky transformation & \texttt{src/innovcal/ca\_rnn/distributions.py} \\
Projected PIT and calibration losses & \texttt{src/innovcal/ca\_rnn/losses.py} \\
Training and early stopping & \texttt{src/innovcal/ca\_rnn/trainer.py} \\
Predictive sampling & \texttt{src/innovcal/ca\_rnn/predict.py} \\
Controlled DGPs & \texttt{src/innovcal/data/simulation.py} \\
Sample-based evaluation & \texttt{src/innovcal/evaluation/metrics.py} \\
Hierarchical paired inference & \texttt{src/innovcal/evaluation/uncertainty.py} \\
Seed and regime experiments & \texttt{src/innovcal/experiments/reproducibility.py} \\
\bottomrule
\end{tabular}
\end{table}

\begin{thebibliography}{9}

\bibitem{huang2024cabl}
J. Huang, S. Park, and O. Simeone,
``Calibration-Aware Bayesian Learning,''
\emph{arXiv preprint arXiv:2305.07504}, version 2, 2024.

\end{thebibliography}

\end{document}
